\oldpage[IV-4]{308}
que $j:U\to X$ est affine implique que $U^0$ s'identifie à $U$ et $j^0$ à $j$,
autrement dit que $j^0(U^0)=U$; or, cela contredit la conclusion de (21.12.6) qui,
sous l'hypothèse de dimension faite implique que $j^0$ doit être surjectif.
\end{proof}

\begin{env}[21.12.8]
\label{IV.21.12.8}
Soient $A$ un anneau local noethérien $Y=\operatorname{Spec}(A)$, $y$ l'unique point
fermé de $Y$, $Y'=Y\setmin\{y\}$. Considérons la condition suivante :
\begin{itemize}
  \item[\textup{(W)}] \emph{Pour tout ouvert $U$ de $Y$ contenu dans $Y'$, ne
  contenant aucune composante irréductible de $Y'$ et tel que l'immersion canonique
  $U\to Y'$ soit affine, $U$ est lui-même un ouvert affine.}
\end{itemize}

Cette condition est entraînée par la suivante :
\begin{itemize}
  \item[\textup{(W')}] \emph{Pour toute partie fermée $T$ de $Y$ dont toute composante
  irréductible est de codimension $1$ dans $Y$, et tel que pour toute composante
  irréductible $Y_i$ de $Y$, on ait $Y_i\cap T\ne\{y\}$, l'ouvert $U=Y\setmin T$ est affine.}
\end{itemize}

En effet, si (W') est vérifiée et si l'ouvert $U$ de $Y$ vérifie l'hypothèse de (W), il
résulte de (21.12.7) qu'aucune composante irréductible de $T=Y\setmin U$ ne peut être de
codimension $\geq2$; comme par ailleurs $U$ ne contient aucune des composantes
irréductibles $Y_i\setmin\{y\}$ de $Y'$, la condition (W') montre que $U$ est affine.

On notera que la condition (W') se simplifie lorsque $Y$ est \emph{irréductible} et est
alors équivalente à la suivante :
\begin{itemize}
  \item[\textup{(W'')}] \emph{Pour toute partie fermée irréductible $T$ de $Y$, de
  codimension $1$ dans $Y$, $Y\setmin T$ est un ouvert affine.}
\end{itemize}

En effet, il est clair que (W') entraîne (W'') lorsque $Y$ est irréductible, et la
réciproque est vraie en considérant les composantes irréductibles de $T$ et en notant
que l'intersection d'un nombre fini d'ouverts affines est un ouvert affine
$(\mathbf I, 5.5.6)$.
\end{env}

\begin{examples}[21.12.9]
\label{IV.21.12.9}
Si $A$ est un anneau local noethérien \emph{factoriel}, il vérifie (W') (et
\emph{a fortiori} (W)), puisque tout idéal premier de hauteur $1$ est principal
$(\mathbf I, 1.3.6)$. Mais il y a des anneaux locaux noethériens non factoriels
vérifiant (W'), par exemple ceux de dimension $\leq1$: en effet, on a remarqué dans la
démonstration de (21.12.6, (i)) que \emph{tous} les ouverts de $Y$ sont alors affines.
On peut d'autre part prouver en utilisant la théorie de la dualité locale (chap.~III,
3\textsuperscript{e} partie) que tout anneau local noethérien de dimension $2$ vérifie
(W').
\end{examples}

L'intérêt de la condition (W) réside dans le résultat suivant :

\begin{proposition}[21.12.10]
\label{IV.21.12.10}
Soient $X$, $Y$ deux préschémas localement noethériens, $g:X\to Y$ un morphisme,
$y$ un point fermé de $Y$, $Y'=Y\setmin\{y\}$, $X'=g^{-1}(Y')$; supposons que le morphisme
$g':X'\to Y'$ restriction de $g$ soit une immersion ouverte, et que l'anneau local
$\mathcal O_{Y,y}$ vérifie la condition \textup{(W) (21.12.8)}. Alors, pour toute composante
irréductible $Z$ de $X_y=g^{-1}(y)$, \textup{ou bien} $Z$ est de codimension $\leq1$ dans $X$,
\textup{ou bien} son point générique est isolé dans $Z$. Si $Z$ est localement de type fini sur
$\boldsymbol{k}(y)$, la seconde alternative implique que $Z$ est réduit à un seul point.
\end{proposition}

\begin{proof}
La dernière assertion de l'énoncé résulte de ce que $Z$ est alors un préschéma de
Jacobson (10.4.7), et comme l'ensemble des points fermés de $Z$ est alors dense dans
$Z$, le point générique de $Z$ ne peut être isolé dans $Z$ que si $Z$ est réduit à un
seul point.

Supposons qu'il existe dans $X_y$ une composante irréductible $Z$ dont le point
générique $z$ n'est pas isolé dans $Z$, et telle que $\operatorname{codim}(Z,X)\geq2$.
La question étant locale
\oldpage[IV-4]{309}
sur $X$ et $Y$ puisque $z$ est non isolé dans $Z$, on peut, en
remplaçant $X$ par un voisinage ouvert de $z$ dans $X$, supposer $X$ et $Y$ affines,
$X_y=Z$ \emph{irréductible, non réduit à un point} et telle que
$\operatorname{codim}(X_y,X)\geq2$. L'image $g(X\setmin X_y)=U$ est un ouvert de $Y'$
isomorphe à $X\setmin X_y$ par hypothèse; montrons qu'en remplaçant au besoin $X$ par un
voisinage ouvert de $z$ dans $X$, on peut supposer que $U$ ne contient aucune des
composantes irréductibles de $Y'$ dont l'adhérence contient $y$.
En effet, on peut tout d'abord supposer que tous les points maximaux de $X$ sont des
générisations de $z$, l'ensemble de ces points étant fini; donc l'ensemble des points
maximaux de $U$ est l'ensemble des images $y_i$ par $g$ des points maximaux $x_i$ de
$X$ (aucun point maximal de $X$ ne pouvant par hypothèse être contenu dans $X_y$,
puisque $\dim(\mathcal O_{X,z})\geq2$). Posons $X_i=\overline{\{x_i\}}$. Par
hypothèse, on a $z\in X_i$ et puisque $z$ n'est pas isolé dans $X_y$, il existe dans
$X_y$ un point $x\ne z$. Comme $X_i\ne Z$, il existe dans $X_i$ un point $t_i$ tel
que si l'on pose $T_i=\overline{\{t_i\}}$, on ait $\dim(\mathcal O_{T_i,x})=1$
(10.5.9); cela entraîne par hypothèse $t_i\notin Z$, donc $t_i$ n'est pas une
générisation de $z$. Remplaçant $X$ par $X'=X\setmin\bigcup_iT_i$, on voit que l'image par
$g$ de $X'\setmin X'_y$ ne contient pas les images $g(t_i)$, donc ne contient aucune des
composantes irréductibles de $Y'$ dont l'adhérence contient $y$.
Cela étant, comme $X$ et $Y$ sont affines, le morphisme $g:X\to Y$ est affine, donc il
en est de même de la restriction $g':X'\to Y'$; puisque $g'$ est une immersion ouverte,
l'immersion canonique $U\to Y'$ est affine. Posons
$Y_1=\operatorname{Spec}(\mathcal O_{Y,y})$, $Y'_1=Y_1\setmin\{y\}=Y'\cap Y_1$,
$U_1=U\cap Y_1$. Ce qui précède prouve que l'immersion canonique $U_1\to Y'_1$ est
affine, donc $U_1$ est un ouvert affine dans $Y_1$ en vertu de (W), en d'autres termes
l'immersion canonique $U_1\to Y_1$ est affine. Mais puisque cette immersion est de
présentation finie (1.6.2), il résulte de (8.10.5, (viii)) appliqué suivant la méthode de
(8.1.2, \emph{a)}) qu'en restreignant au besoin $Y$ à un voisinage ouvert affine de $y$, on peut
supposer que l'immersion $U\to Y$ est affine. On en conclurait que l'ouvert $X\setmin X_y$ de
$X$, isomorphe à $U$, serait affine, ce qui contredirait (21.12.7); la proposition est
ainsi démontrée.
\end{proof}

\begin{corollary}[21.12.11]
\label{IV.21.12.11}
Soient $X$, $Y$ deux préschémas localement noethériens, $X$ étant supposé
irréductible; soit $g:X\to Y$ un morphisme localement de type fini, et soit $V$ le plus
grand ouvert de $X$ tel que la restriction $g|V:V\to Y$ soit un isomorphisme local.
Supposons que pour tout point $y\in Y$, l'anneau local $\mathcal O_{Y,y}$ vérifie la
condition \textup{(W) (21.12.8)}. Alors toute composante irréductible de $T=X\setmin V$ est, \textup{ou bien}
de codimension $\leq1$, \textup{ou bien} telle que son point générique $z$ soit isolé dans
$g^{-1}(g(z))$.
\end{corollary}

\begin{proof}
Posons $g(z)=y$. La question ne dépendant que de la fibre $g^{-1}(y)$ et de l'anneau
$\mathcal O_{X,z}$, on peut par changement de base, se borner au cas où
$Y=\operatorname{Spec}(\mathcal O_{Y,y})$ et où $X$ est affine
$((\mathbf I, 3.6.5)$ et $(\mathbf I, 4.5.5))$; remplaçant éventuellement $X$ par un
voisinage ouvert de $z$ dans $X$, on peut supposer que $T$ est irréductible; en outre on
peut se borner au cas où $V$ est non vide. Le morphisme $g$ peut donc être supposé
\emph{séparé} et $y$ fermé dans $Y$; puisque la restriction $g|V:V\to Y$ est un isomorphisme
local et que l'ouvert non vide $V$ de $X$ est irréductible, il résulte de
$(\mathbf I, 8.2.8)$ que $g|V:V\to Y$ est une \emph{immersion ouverte}. La restriction
$g|(V\cap X_y):V\cap X_y\to\{y\}=\operatorname{Spec}(\boldsymbol{k}(y))$ est donc aussi une
immersion ouverte, ce qui montre que $V\cap X_y$ est, soit vide, soit réduit à un point
$x$ rationnel
